\item La Ley de Boyle establece que cuando una muestra de gas se comprime a temperatura constante, la presión $P$ y el volumen $V$ satisfacen la ecuación $PV = C$, en donde $C$ es una constante. Supóngase que en cierto instante el volumen es de $600 \text{ cm}^3$, la presión de $150 \text{ kPa}$ y que ésta aumenta a razón de $20 \text{ kPa/min}$. ¿A qué velocidad disminuye el volumen en dicho instante?
